\documentclass[unnumsec,webpdf,contemporary,large]{oup-authoring-template}
% The page headers draw a society logo here; leave it empty unless the journal gives one.
\providecommand{\societylogo}{}

\begin{document}

\journaltitle{Journal Title}
\DOI{DOI}
\copyrightyear{2026}
\pubyear{2026}
\access{Advance Access Publication Date: Day Month Year}
\appnotes{Original Article}

\firstpage{1}

\title[Short title]{Title of the Study: Outcomes After an Intervention in Adults With a Chronic Condition}

\author[1,$\ast$]{First Author\ORCID{0000-0000-0000-0000}}
\author[2]{Second Author}
\author[1,2]{Third Author}


\address[1]{\orgdiv{Department of Medicine}, \orgname{University Hospital}, \orgaddress{\street{Street}, \postcode{00000}, \state{City}, \country{Country}}}
\address[2]{\orgdiv{School of Public Health}, \orgname{University}, \orgaddress{\state{City}, \country{Country}}}

\corresp[$\ast$]{Corresponding author. \href{mailto:first.author@example.org}{first.author@example.org}}

\received{Date}{0}{Year}
\revised{Date}{0}{Year}
\accepted{Date}{0}{Year}

\abstract{\textbf{Background:} State why the question matters and what is unknown.
\textbf{Methods:} Describe the design, setting, participants and main outcome.
\textbf{Results:} Give the main findings with effect sizes and 95\% confidence intervals.
\textbf{Conclusions:} Say what the findings mean for practice or research.}

\keywords{randomized controlled trial, chronic disease, quality of life}

\maketitle

\section{Introduction}
Summarise what is known and the gap this study addresses. The trial is
reported following CONSORT~\citep{schulz2010consort} and was conducted in line
with the Declaration of Helsinki~\citep{wma2013helsinki}.

\section{Methods}
\subsection{Trial design and participants}
Describe the design, setting, dates, eligibility criteria and randomization.

\subsection{Outcomes}
Define the primary and secondary outcomes and how they were measured.

\subsection{Statistical analysis}
Describe the sample size calculation and the models used. The treatment effect
on the primary outcome was estimated as
\begin{equation}
  y_i = \beta_0 + \beta_1\,\mathrm{treatment}_i + \beta_2\,\mathrm{baseline}_i + \varepsilon_i .
\end{equation}

\section{Results}
Table~\ref{tab:outcomes} gives the outcomes in each group, and
Figure~\ref{fig:flow} the flow of participants.

\begin{table}[!t]
\caption{Primary and secondary outcomes at 12 months.\label{tab:outcomes}}
\begin{tabular}{lrrr}
\toprule
Outcome & Intervention & Control & Difference (95\% CI) \\
\midrule
Primary score & 62.1 & 55.4 & 6.7 (3.2 to 10.2) \\
Hospital admission, \% & 8.2 & 11.9 & $-$3.7 ($-$7.5 to 0.1) \\
\botrule
\end{tabular}
\begin{tablenotes}%
\item CI, confidence interval.
\end{tablenotes}
\end{table}

\begin{figure}[!t]
\centering
\includegraphics[width=0.6\linewidth]{example-image}
\caption{Flow of participants through the trial. Upload an image and use its file name here.}\label{fig:flow}
\end{figure}

\section{Discussion}
Interpret the findings, compare them with earlier studies, and discuss
strengths and limitations.

\section{Conclusion}
State the main conclusion in a few sentences.

\section{Trial registration}
Give the registry and the registration number.

\section{Funding}
List the funders and grant numbers, or state that there was no specific funding.

\section{Conflict of interest}
The authors declare no conflicts of interest.

\section{Data availability}
Say where the data can be obtained, and under which conditions.

\bibliographystyle{oup-abbrvnat}
\bibliography{references}

\end{document}
